\section{Data and Methods}
\subsection{Data, Targets and Quality}
OpenAQ preserves upstream CPCB-provider station identity; its measurements are
not an official CPCB bulletin or licensed CCR export \cite{openaq2026}.
The inherited catalogue contains 481 eligible stations across 255 parsed cities;
the metadata-only stratified sample selects 181 stations in 100 cities. Current
catalogue selection can introduce survivorship bias and is disclosed.
The locked v3.1 window is 2025-02-19 through 2026-10-07. The protocol was
locked at 2026-10-11T18:32:05.939254+00:00; its prior version and hash remain in the lock history.
Its immutable hash is recorded in the result artifact. This acquisition checkpoint has
179 usable stations, 100 cities, 64799 valid targets and 42481
retained rejected station-days. These are partial-collection counts, not
national study acceptance.

The measured semantic dataset and preliminary input share SHA-256 prefix \texttt{4761a5da60dde059}; the immutable preliminary run has prefix \texttt{287c87db04a6b294}. Full hashes, denominators and artifact locations are recorded in the claim-evidence manifest.

The single CPCB-style engine applies pollutant breakpoints, at least sixteen
distinct hours per pollutant, at least three supported pollutants and a required
particulate species \cite{cpcb2014aqi}. Particulates and the designated gases use
24-hour means; CO and ozone use the maximum complete trailing eight-hour mean.
Targets are labelled AirGuard CPCB-style AQI. Exact zeros, provider flags, poor
coverage and six-hour identical runs are rejected. PM2.5 above PM10 and potential
spikes are flagged and retained. Raw caches and row-addressable rejection ledgers
preserve exclusions. Unsupported ppb conversion is refused.

Severe AQI is an open lower bound of at least 401. The checkpoint contains
383 censored Severe days. They cannot supply exact residual labels or
identified MAE, RMSE, bias or coefficient-of-determination values. Category,
exceedance, forecasts below 401 and lower-bound shortfall remain estimable.
% BEGIN TARGET_METHOD
\subsubsection{Concentration-to-index transformation}
Within a bounded CPCB interval $[C_L,C_H]$ with index endpoints $[I_L,I_H]$,
the pollutant sub-index is
\begin{equation}
I_p(C)=I_L+\frac{I_H-I_L}{C_H-C_L}(C-C_L),\quad I=\max_{p\in S}I_p.
\end{equation}
The daily support $S$ requires at least three sufficiently covered species,
including particulate matter. It is recorded as an outcome annotation, never
used to select predictors retrospectively. In the current valid development
join, $S$ equals the prerequest frozen PM2.5/PM10/O3 support. Rejected rows
retain reasons and provenance rather than disappearing from denominator
counts. An open concentration interval does not justify linear extrapolation
to a fictitious exact Severe target.
\subsubsection{Quality-control interpretation}
The six-hour stuck-run rule rejects the complete identical sequence, not only
its final bin. Exact-zero rejection cannot distinguish a genuinely tiny
ozone concentration from a telemetry zero and is therefore a disclosed
preregistered trade-off. Conversely, automatic spike deletion could erase
real pollution episodes. PM2.5/PM10 inconsistencies and potential spikes
remain flagged observations. Sensor rejection, unavailable target and
model error are distinct outcomes; their counts are reported separately.
% END TARGET_METHOD

\subsection{Issue-Time Availability}
Issue time is 20:00 IST on the preceding day. Observation bins must end by
17:00 IST; daily history is D-2 or earlier. Historical OpenAQ publication time
is unknown, so measured-history analyses explicitly assume a healthy feed.
Archived weather uses the documented 48-hour Previous Runs offset
at \url{https://open-meteo.com/en/docs/previous-runs-api}. The Air Quality API documentation is a separate
source \cite{openmeteo2026}. FIRMS acquisition time alone does not establish release
time. CAMS past fields without their producing run are excluded from the primary
experiment. Every admitted snapshot passes its guard.


\subsection{Primary Forecast and Correction}
The genuine CAMS global forecast archive \cite{cams2024forecasts} must retain
initialisation, lead, valid time, units and evidence of pre-issue availability.
Reanalysis is prohibited as a replacement. The locked v3.1 protocol specifies
target-support-matched CAMS scoring and a six-pollutant secondary comparator.
Later target support is an evaluation annotation and cannot become a predictor.
The partial join contains 7,560 station-days, 5,506 valid targets and 2,054 retained rejected targets across 42 dates. All 7,560 rows pass the primary availability guard. Every valid target in this subset uses the frozen PM2.5/PM10/O3 support; outcome support remains an evaluation annotation and is excluded from predictors. Full feature integration and the primary runner remain incomplete. The locked comparison is preserved. No primary model has been fitted and no final-test outcomes have been accessed.
All methods share observations and evaluation windows. With frozen CAMS index
$b_{s,d}$ and exact measured index $y_{s,d}$, residual learning uses
\begin{equation}r_{s,d}=y_{s,d}-b_{s,d},\qquad
\widehat y_{s,d}=b_{s,d}+\widehat r_{s,d}.\end{equation}
The ladder comprises raw CAMS, direct ML, residual Ridge, residual LightGBM and
residual XGBoost, with persistence and prior-only climatology as operational
floors. There is no assumption that correction wins. The legacy city model has
a different period, target and input representation, and is reported separately.

\begin{table}[h]\centering\caption{Primary comparison contract}
\begin{tabular}{ll}\toprule Method & Result \\\midrule
Raw operational CAMS & Pending archive \\
Direct ML & Pending archive \\
Residual Ridge & Pending archive \\
Residual LightGBM & Pending archive \\
Residual XGBoost & Pending archive \\\bottomrule\end{tabular}\end{table}






% BEGIN CURRENT CAMS EVIDENCE
\subsubsection{Verified operational archive checkpoint}
The single accepted probe is a genuine forecast (GRIB type fc), initialized
2025-02-18 00:00 UTC with a 24-hour lead and valid at 2025-02-19 00:00 UTC.
Its 134-byte PM2.5 field contains nine finite values on a 0.4-degree grid.
The development manifest currently verifies 6 of 36 chunks.
Requests cover only the frozen station bounding box and development dates.
Initialization plus lead reconstructs valid time in every admitted field.
Surface PM mass concentrations use $10^9$ to convert kg m$^{-3}$ to
$\mu$g m$^{-3}$. Gas mixing ratios use the same-run lowest-level
temperature $T$, specific humidity $q$ and hybrid pressure $p$:
\begin{equation}
\rho=\frac{p}{287.05T(1+0.6078q)},\quad C_{\mu g/m^3}=10^9\rho x.
\end{equation}
CO instead uses $10^6$ for mg m$^{-3}$. Model-level fields are linearly
interpolated within the same producing run before density conversion;
no extrapolation, reanalysis or post-issue run is admitted. Seven preceding
hourly bins support trailing eight-hour CO/O3 windows.
Historical per-run delivery logs are unavailable: pre-issue availability
uses the documented nominal dissemination bound, explicitly an assumption.
The later ADS retrieval timestamp is provenance, not historical publication.
% END CURRENT CAMS EVIDENCE


\subsection{Selection and Ablation}
Model settings use training-only inner validation. The locked protocol fixes
90 training days initially, 30 calibration days, 30 test days, a 30-day step and
one-day embargo gaps, with expanding training windows. Each fold is reported.
The latest 90 dates, 2026-07-10 through 2026-10-07, are sealed.
The referenced artifact declares no final-test scoring. The primary selection rule is lowest mean rolling-fold MAE, breaking ties toward Ridge and residual formulations; raw CAMS is retained if no learned method beats it. A committed selection lock, complete isolation audit and fresh user approval are required before final evaluation. Feature ablation separates CAMS, measured history,
forecast weather, calendar, spatial context and optional fires. Fire admission
requires a date-block interval supporting its benefit. An additional sliding
90-day non-CAMS sensitivity is labelled separately.
[RESULT PENDING CAMS ARCHIVE EXPERIMENT]


\subsection{Temporal and Spatial Verification}
Primary evaluation will report exact-target MAE, RMSE, signed forecast bias,
coefficient of determination, six-category diagnostics, threshold detection,
false alarms and censored-Severe shortfall. Poor, Very Poor and Severe events
receive explicit diagnostics. Seasonal and regional groups below thirty pairs
are listed with counts and withheld metrics.
Spatial tests isolate cities, 25-km connected urban clusters, stations and
regions. Local-history and monitor-free inference are distinct conditions.
Held-out station targets cannot determine training eligibility. The all-station
reference is scored on the identical held-out test rows.

Paired comparison uses seven-day moving date blocks with 5,000 resamples and
a two-sided Diebold--Mariano daily loss comparison with Newey--West variance
and Holm multiplicity adjustment \cite{diebold1995comparing}. A confidence
interval containing zero does not establish improvement.

\subsection{Uncertainty Calibration}
Split conformal intervals at 80\% and 90\%, quantile calibration and
forecast-category Mondrian calibration use only the preceding calibration
window. Coverage and widths are reported by category, season and region.
For a censored Severe outcome, an upper bound reaching 401 is merely compatible
with the known lower bound; it does not identify true-value coverage.
Calibration under temporal and spatial shift is evaluated empirically;
exchangeability guarantees are not asserted for dependent station-days
\cite{gibbs2021aci,xu2021enbpi}. Proper uncertainty scoring requires identified
outcomes \cite{gneiting2007scoring}.





% BEGIN PRELIMINARY COVERAGE
\begin{table}[h]\centering\caption{NON-CAMS preliminary first-fold calibration failures}
\begin{tabular}{lrr}\toprule Method & Exact coverage & Mean width\\\midrule
80\% split & 0.855 & 95.5\\
90\% split & 0.944 & 120.6\\
\bottomrule\end{tabular}
\end{table}
Marginal coverage conceals important failures: neither interval covered any
of the 35 exact Very Poor outcomes in this first fold. Poor-or-worse coverage
was also zero on 48 exact pairs. These are preliminary validation failures,
not primary CAMS or final-test findings; they preclude claiming reliable
high-pollution uncertainty from average coverage alone.
% END PRELIMINARY COVERAGE


\subsection{Vulnerability-Aware Advisory Evaluation}
Declared exposure strata include general adults, children, older adults,
asthma/COPD, heart conditions, pregnancy and outdoor workers. They are not
personal medical-risk predictions. P0 uses the point forecast; P1 the calibrated
median; P2 the calibrated 0.75 quantile; P3 the smallest calibration-window
candidate achieving the corrected empirical under-warning aim $(k+1)/(n+1)$.
The aims are 10\% for general adults and 5\% for the other declared strata.
If no candidate qualifies, the policy reports that fact.
Under-warning, over-warning, serious under-warning and high-pollution misses
are reported alongside date-block confidence intervals. A lower warning rate
does not demonstrate a health benefit. Unsupported cohort coefficients remain
unadopted. [RESULT PENDING CAMS ARCHIVE EXPERIMENT]
% BEGIN ADVISORY_METHOD
\subsubsection{Policy losses and calibration boundary}
For ordered observed category $c$ and warning category $w$, the basic losses
are $L_-=\mathbf{1}\{w<c\}$ and $L_+=\mathbf{1}\{w>c\}$.
Serious under-warning, missed Poor-or-worse and missed Very-Poor-or-worse
events remain separate diagnostics. Lower under-warning can require more
false alarms; neither objective is optimized on a test window.
The candidate quantiles are fixed before evaluation. P3 chooses on the
calibration window only and then reports its frozen test-window behavior,
including failure to find any feasible candidate. Its corrected empirical
rate $(k+1)/(n+1)$ is a policy selection statistic, not an individual disease
probability or a proven dependent-data risk bound.
\subsubsection{Finite-sample intervals and failure reporting}
For exact calibration residuals $e_i=|y_i-\widehat y_i|$, split conformal
uses the order statistic with rank $\lceil(n+1)(1-\alpha)\rceil$.
The interval is $[\max(0,\widehat y-q),\widehat y+q]$; missing or inadequate
calibration support is reported rather than assigned a radius of zero.
Mondrian groups use the forecast category, so observed test categories cannot
choose an interval. Quantile crossing is repaired using predictions only.
Calibration labels exclude censored Severe values from exact residuals;
Severe category/threshold behavior is still reported. Under temporal/spatial
dependence, coverage is an empirical finding, not automatically a theorem.
\subsubsection{Exposure guidance and clinical scope}
Declared vulnerability strata set environmental-warning tolerances. Shared
forecast outcomes do not constitute separate clinical cohort studies. The
system does not predict diagnoses, prescribe medication, estimate personal
health probabilities or demonstrate avoided illness. Its interface labels
measurement versus modelling, observation age, unsupported intervals and
raw-model uncertainty. Clinical coefficients remain unvalidated.
% END ADVISORY_METHOD

\begin{table}[t]\centering\caption{Dataset audit. Full measured counts are inherited descriptive metadata; only development rows enter alignment.}
\begin{tabular}{lr}\toprule Item & Count\\\midrule
Measured stations & 180\\
Usable stations & 179\\
Usable cities & 100\\
Retained station-days & 107,280\\
Valid targets & 64,799\\
Rejected targets & 42,481\\
Censored Severe targets & 383\\
\midrule
Aligned development rows & 7,560\\
Aligned dates & 42\\
Aligned valid targets & 5,506\\
Aligned rejected targets & 2,054\\
Aligned Severe targets & 2\\
\bottomrule\end{tabular}\end{table}
\begin{table}[t]\centering\caption{Preregistered search space, not selected hyperparameters. All model selection remains pending.}
\begin{tabular}{lp{.57\columnwidth}}\toprule Candidate & Training-only choices\\\midrule
Ridge & $\alpha\in\{0.1,1,10,100\}$\\
LightGBM & Leaves 15/31; minimum child samples 20/60\\
XGBoost & Depth 4/6; minimum child weight 5/20\\
Both boosters & Learning rate 0.05; maximum 600 rounds; early stopping 40 rounds; subsample and column fraction 0.8\\
All learners & Seed 42; direct and residual formulations\\
Selection & Mean rolling-fold MAE; simpler learner and residual formulation break ties\\
\bottomrule\end{tabular}\end{table}

\subsubsection{Baseline estimation and operational missingness}
Climatology is estimated from exact training targets in a station's calendar
month, with station and global means as explicit fallbacks. These statistics
are re-estimated inside each training fold. Persistence uses the measured
rolling 24-hour index ending at the cutoff, rather than the not-yet-complete
day immediately before the target. Consequently a persistence forecast can
be missing even when the eventual target is valid. The primary comparison
requires exact targets, the CAMS baseline and persistence jointly; the
secondary set relaxes only the persistence requirement. Reporting these
denominators prevents a method gaining an advantage by being evaluated on
easier, more complete days.

The implementation joins on station identifier and calendar date and rejects
duplicate keys or mismatched station coordinates. Calendar lags are looked up
by date, rather than shifted across whichever observations happen to be
present. An absent day is therefore missing, not replaced by an older reading
with the wrong lead. Censored historical values retain a lower-bound indicator;
they do not supply fabricated exact residual-bias labels. Numerical missingness
is handled by training-only imputation with indicators. This choice describes
the current implementation and remains subject to the development audit;
no imputer is fitted to a validation or final window.

Nearby-station summaries use only donor rolling measurements whose bins
ended by the cutoff. Self-donations are excluded. Spatial inference without
a local monitor must mask unavailable donor stations as well as the held-out
station's own history. A held-out target is never a contemporaneous neighbour
feature. Elevation is static cached terrain metadata, not an outcome-derived
station-selection criterion. The acquired surface cube also provides maxima,
ranges, standard deviations and daytime-minus-nighttime concentrations; all
come from the same producing run as the daily CAMS baseline.

\subsubsection{Scope of the operational reconstruction}
The locked YAML preserves the earlier proxy-source disclosure and its hash;
it has not been silently rewritten to claim historical operational access.
The separately frozen acquisition manifest identifies the genuine ADS
forecast stream now being integrated. That provenance improvement does not
relax the original QC, date splits, censoring or selection criteria. It also
does not prove each forecast's historical delivery time. Nominal publication
and healthy-feed assumptions remain explicit limitations of retrospective
reconstruction, to be checked prospectively.
